\section{Discussion}
\label{sec:discussion}
The observations of Section~\ref{sec:results} come from 68 real queries, are benchmark-derived and are partly clustered. We use them to say where the design is limited, and not to estimate how accurate it is.

\textbf{Coverage is limited by competition, not by weak evidence.} The top candidate's score was high in almost every case (median 0.88), yet only 4 of 50 smoke queries and none of 18 exploratory queries were assigned automatically under C1. The reason is visible in the gate counts: a margin below 0.15 or a dominance ratio below 1.35 in the large majority of cases (37 and 44 of 50), and fewer than two observed signals in 14 of 50. In a block of same-name authors, several candidates share the printed name and often a topic and affiliation, so each scores close to the others. A fixed-weight score over name, institution and years separates people poorly when the name signal, which carries the largest weight, is saturated for all of them. Raising coverage by lowering the margin or dominance thresholds would convert exactly these cases into automatic assignments of a candidate that is barely distinguished from its competitor, so the low coverage is the policy doing what it was designed to do. Whether it does so at the right price is the question the verified evaluation must answer.

\textbf{The heuristic ordering did not beat the provider's order.} Under C1 the top candidate by HAMMER's score agreed with the benchmark label in 34\% and 33\% of the queries, whereas the provider's first candidate agreed in 44\% and 39\%. With so few cases the difference is not established, but it removes any basis for claiming that the five-signal score improves ranking, and it suggests that candidates are being reordered mostly by name similarity, which the provider has already used. We do not know how far this reflects the weights, the missing coauthor signal on this benchmark, or the limited context supplied; the weights were never tuned and the score is not a probability.

\textbf{Retrieval is a ceiling.} The profile carrying the benchmark ORCID was missing from the pool for about a quarter of the queries (12 of 50, 5 of 18), and no ranking or decision rule can recover it. Some of these are retrieval failures, for example profiles split across several identifiers or a name form that the provider's search does not match, and some are label artefacts, such as a right-person profile without an ORCID. Telling them apart needs the independent review, which is why the design reports retrieval and assignment separately and keeps the verified retrieval measures apart from the benchmark-derived ones.

\textbf{The safeguards matter where evidence is cheap.} Under C1 the safeguards B1 to B5 changed no outcome, but under the leakage-prone C1x they withheld automation in 19 of 26 smoke cases and 5 of 8 exploratory cases. This is consistent with their purpose, since publication overlap with a group member's own paper is an easy and unreliable signal, but whether the withheld cases were wrong is unknown. We do not claim that the safeguards improve precision.

\textbf{What the label cannot say.} One of the four automatic smoke assignments disagreed with the benchmark. It may be an error of the system or of the label, and only independent review can decide. This single case shows why the design verifies every automatic assignment, and why the benchmark agreement reported here is not called precision.

\textbf{Consequences for a verified evaluation.} At the observed automatic rate of about 6\% a verified evaluation targeting 32 automatic assignments needs several hundred to about a thousand primary cases (Table~\ref{tab:exp-planning}), and reviewers would have to judge every automatic assignment and an audit sample. The cost of verification therefore scales with the coverage, which is itself small. Raising coverage in a defensible way, for example by supplying richer context or by resolving institutions to identifiers before scoring, would reduce the review burden before it improves any precision figure.

\textbf{Limits of these observations.} The samples are small (68 queries), the smoke sample is clustered into 23 blocks, the exploratory sample is a random subset of a stratified draw with no L or XL case, the label is provisional, a single provider and date are involved, and several comparisons rest on one to four cases. They are hypotheses about the behaviour of the approach, to be confirmed or rejected by the planned evaluation.
